\subsection{محاسبه وزن‌های محلی در مدیر خوشه}
\label{sec:cluster-manager}

مدیر خوشه برای هر مبدأ $e$ و سرویس $s$، وزن نمونه‌های داخلی را جدا محاسبه می‌کند. بنابراین، دو مبدأ می‌توانند درخواست‌های یک سرویس را با نسبت‌های متفاوت بین مقصدها تقسیم کنند. محل این محاسبه، بخش «هزینه محلی» در شکل~\ref{fig:algorithm-architecture} است.

در این بخش، $\ell_{es}$ تعداد درخواست بر ثانیه‌ای است که باید داخل خوشه پاسخ بگیرد؛ قاعده تعیین آن در بخش~\ref{sec:integration} می‌آید. ظرفیت هر مقصد به کانتینر همان سرویس تعلق دارد و در طول هر حل ثابت است.

\subsubsection{برآورد خطی ظرفیت از مصرف پردازنده}

برای هر کانتینر، سهم پردازنده $g_{si}$، زمان مصرف آن $T^{\mathrm{CPU}}_{si}$ و شمار پردازش‌های پایان‌یافته $N^{\mathrm{proc}}_{si}$ از آخرین پنجره معتبر دریافت می‌شوند. طول پیش‌فرض پنجره $\Delta=5$ ثانیه است. زمان مصرف، اختلاف شمارنده تجمعی پردازنده در آغاز و پایان همان پنجره است.

میانگین هسته مصرف‌شده، نسبت استفاده از سهم اختصاص‌یافته و نرخ پردازش مشاهده‌شده برابرند با
\begin{align}
 \overline c_{si}&=\frac{T^{\mathrm{CPU}}_{si}}{\Delta},\qquad
 u_{si}=\frac{\overline c_{si}}{g_{si}},
 \label{eq:container-utilization}\\
 y_{si}&=\frac{N^{\mathrm{proc}}_{si}}{\Delta}.
 \label{eq:processed-throughput}
\end{align}
با فرض تناسب خطی نرخ پردازش با هسته مصرف‌شده، ضریب خطی و ظرفیت کانتینر چنین به دست می‌آیند:
\begin{equation}
 \gamma_{si}=\frac{y_{si}}{\overline c_{si}}
            =\frac{N^{\mathrm{proc}}_{si}}{T^{\mathrm{CPU}}_{si}},\qquad
 \mu_{si}=\gamma_{si}g_{si}=\frac{y_{si}}{u_{si}}.
 \label{eq:cpu-linear-capacity}
\end{equation}
ضریب $\gamma_{si}$ تعداد درخواست بر ثانیه به ازای یک هسته مصرف‌شده را بیان می‌کند. مقدار $\mu_{si}$ نرخ برآوردشده در استفاده کامل از سهم پردازنده کانتینر است. اگر مصرف به صورت درصد $p_{si}$ گزارش شود، در این رابطه $u_{si}=p_{si}/100$ است.

این مدل برای ترکیب درخواست نسبتا ثابت و کار وابسته به پردازنده مناسب است. $\mu_{si}$ برآورد ظرفیت است، نه بیشینه تضمین‌شده؛ تغییر نوع درخواست یا سهم پردازنده می‌تواند آن را تغییر دهد.

اگر پردازش معتبر وجود نداشته باشد یا زمان مصرف صفر باشد، نسبت ساخته نمی‌شود. آخرین برآورد معتبر همان پیکربندی یا اندازه‌گیری راه‌اندازی به کار می‌رود. نمونه بدون برآورد معتبر در ظرفیت عادی خوشه شمرده نمی‌شود.

\subsubsection{شمار درخواست از همه مبدأها}

شمار پردازش پایان‌یافته در مقصد، همه فرستنده‌ها را پوشش می‌دهد:
\begin{equation}
 N^{\mathrm{proc}}_{si}=\sum_qN^{\mathrm{proc}}_{qsi},\qquad
 N^{\mathrm{in}}_{si}=\sum_qN^{\mathrm{in}}_{qsi},\qquad
 \lambda^{\mathrm{obs}}_{si}=\frac{N^{\mathrm{in}}_{si}}{\Delta}.
 \label{eq:all-source-observation}
\end{equation}
شاخص $q$ مبدأهای محلی و خارجی را شامل می‌شود. اگر شمارنده مقصد از ابتدا کل ورودی یا پردازش را ثبت کند، به دریافت تفکیک مبدأ نیاز نیست و جمع دوباره انجام نمی‌شود. هر تلاش در هر شمارنده مقصد فقط یک بار ثبت می‌شود.

در $N^{\mathrm{proc}}_{si}$، پردازش‌های موفق و ناموفق انجام‌شده شمرده می‌شوند؛ درخواستی که پیش از اجرا رد شده است، پردازش انجام‌شده محسوب نمی‌شود. شمار ورود ممکن است درخواست منتظر صف را هم داشته باشد؛ بنابراین در برآورد $\mu_{si}$ جای شمار پردازش پایان‌یافته را نمی‌گیرد.

نرخ $\lambda^{\mathrm{obs}}_{si}$ برای پایش و مرجع صف است و از ظرفیت کم نمی‌شود. ظرفیت از پردازش پنجره قبل برآورد می‌شود؛ حل وزن‌ها تخصیص تازه را می‌سازد و بار قبلی دوباره وارد تابع نمی‌شود.

\subsubsection{وزن‌ها، حاشیه اطمینان و سقف مشترک مبدأها}

مجموعه $\mathcal A_{es}=\mathcal A_{as}$، نمونه‌های مجاز، سالم و آماده سرویس در خوشه مبدأ با برآورد ظرفیت معتبر، مثبت و متناهی است. برای هر سطر $(e,s)$ با نرخ محلی مثبت داریم
\begin{align}
 \sum_{i\in\mathcal A_{es}}w_{esi}=1,\qquad w_{esi}&\geq0,
 \label{eq:row-weights}\\
 r_{esi}=\ell_{es}w_{esi},\qquad
 \sum_{i\in\mathcal A_{es}}r_{esi}&=\ell_{es}.
 \label{eq:load-weight}
\end{align}
هر مبدأ سطر وزن مخصوص خود را دارد؛ سقف ظرفیت مقصدها میان این سطرها مشترک است. اگر نرخ محلی صفر باشد، حل آن سطر انجام نمی‌شود و نرخ‌های محلی آن صفرند. مجموع وزن مبدأهای مختلف برای یک مقصد لازم نیست یک باشد.

با حاشیه اطمینان $0<\varepsilon<1$، سقف هر کانتینر و ظرفیت محلی سرویس عبارت‌اند از
\begin{align}
 U_{si}&=(1-\varepsilon)\mu_{si},
 \label{eq:global-cap}\\
 K_{as}&=\sum_{i\in\mathcal A_{as}}U_{si}.
 \label{eq:local-capacity}
\end{align}
هر کانتینر با سهم پردازنده خودش شمرده می‌شود؛ نرخ سرویس‌های دیگر وارد ظرفیت سرویس جاری نمی‌شود.

برای آنکه چند مبدأ ظرفیت یک کانتینر را چند بار مصرف نکنند، مجموع نرخ‌های تازه تخصیص‌یافته با $R^{\mathrm{plan}}_{si}$ نگهداری می‌شود. در آغاز حل هر سرویس، این مقدار برای همه مقصدها صفر است. مبدأها با ترتیب ثابت حل می‌شوند و سقف سطر جاری برابر است با
\begin{equation}
 R^{\mathrm{max}}_{esi}=\max\{0,U_{si}-R^{\mathrm{plan}}_{si}\}.
 \label{eq:remaining}
\end{equation}
این جمع فقط تخصیص تازه همان سرویس در همین چرخه است؛ $\mu_{si}$ در مخرج هزینه ثابت می‌ماند.

مجموعه $\cI_{es}$ مقصدهای با سقف مثبت را نشان می‌دهد. نمونه‌های بیرون از این مجموعه در حل عادی نرخ صفر دارند و برای سایر نمونه‌ها داریم
\begin{equation}
 0\leq r_{esi}\leq R^{\mathrm{max}}_{esi}\leq U_{si}<\mu_{si}.
 \label{eq:domain}
\end{equation}
ازاین‌رو، مخرج هزینه دست‌کم $\varepsilon\mu_{si}$ مثبت می‌ماند. شرط پذیرش محلی $\sum_e\ell_{es}\leq K_{as}$ در بخش ترکیب رعایت می‌شود. با مجموعه مقصد مشترک، ظرفیت باقی‌مانده پس از هر سطر برای نرخ‌های محلی حل‌نشده کافی است.

اگر $q\prec e$ مبدأهای حل‌شده پیش از مبدأ $e$ را نشان دهد، کنترل امکان تخصیص هر سطر چنین نوشته می‌شود:
\begin{equation}
 \sum_i R^{\mathrm{max}}_{esi}
   =K_{as}-\sum_{q\prec e}\ell_{qs}\geq\ell_{es}.
 \label{eq:row-feasibility}
\end{equation}
پس رزرو سطرهای قبلی، امکان تخصیص نرخ محلی سطرهای بعدی را از بین نمی‌برد.


\subsubsection{هزینه ثابت مسیر و مشاهده صف کانتینر}

هزینه ثابت هر مسیر از میانگین پینگ، جریمه شکست و مازاد انتظار صف ساخته می‌شود:
\begin{equation}
 d^{\mathrm{fix}}_{esi}=\overline d^{\mathrm{net}}_{ei}
                   +P_s\widehat f_{esi}+\delta^Q_{si}.
 \label{eq:fixed-cost}
\end{equation}
پارامتر $P_s\geq0$ جریمه تعیین‌شده برای هر شکست، بر حسب ثانیه است. پینگ و نتیجه درخواست در مبدأ ثبت می‌شوند. انتظار صف در کانتینر مقصد اندازه‌گیری می‌شود و به درج زمان پردازش در پاسخ سرویس نیاز ندارد.

برای $n_{si}^{\mathrm{start}}>0$ شروع خدمت در کانتینر، میانگین انتظار مشاهده‌شده چنین است:
\begin{equation}
 \overline W_{si}^Q=\frac{1}{n_{si}^{\mathrm{start}}}
 \sum_{j=1}^{n_{si}^{\mathrm{start}}}
 (t_j^{\mathrm{start}}-t_j^{\mathrm{enqueue}}).
 \label{eq:queue-measurement}
\end{equation}
مرجع ساده صف از مدل $M/M/1$ محاسبه می‌شود. برای حفظ مؤلفه صف روش پایه، فقط مازاد انتظار نسبت به این مرجع به هزینه ثابت افزوده می‌شود:
\begin{align}
 \lambda_{si}^{\mathrm{ref}}&=\min\{\lambda_{si}^{\mathrm{obs}},U_{si}\},\\
 W_{si}^{Q,\mathrm{ref}}&=
 \frac{\lambda_{si}^{\mathrm{ref}}}
      {\mu_{si}(\mu_{si}-\lambda_{si}^{\mathrm{ref}})},\\
 \delta^Q_{si}&=\max\{0,\overline W_{si}^Q-W_{si}^{Q,\mathrm{ref}}\}.
 \label{eq:queue-residual}
\end{align}
نرخ مرجع از کل ورودی کانتینر، شامل تمام مبدأها، ساخته می‌شود. این مرجع برای کانتینر چندرشته‌ای تقریب است؛ مازاد صف، اصلاح تجربی ثابت این چرخه است. اگر داده انتظار موجود نباشد، مؤلفه با پرچم «ناموجود» کنار گذاشته می‌شود؛ نبود داده به معنی صف اندازه‌گیری‌شده صفر نیست.

\Needspace{10\baselineskip}
\begin{longtable}{p{3.1cm}p{8.1cm}p{2.4cm}}
\caption{پارامترهای ظرفیت خطی کانتینر و تخصیص محلی}\label{tab:symbols}\\
\toprule نماد & معنا و روش تعیین & واحد \\
\midrule\endfirsthead
\multicolumn{3}{c}{ادامه جدول \ref{tab:symbols}}\\
\toprule نماد & معنا و روش تعیین & واحد \\
\midrule\endhead\bottomrule\endfoot
$g_{si}$ & سهم قابل اتکای پردازنده همان کانتینر. & هسته \\
$T^{\mathrm{CPU}}_{si}$ & زمان مصرف پردازنده در پنجره قبل. & هسته‌ثانیه \\
$N^{\mathrm{proc}}_{si}$ & کل پردازش‌های پایان‌یافته از همه مبدأها. & درخواست \\
$u_{si},y_{si}$ & نسبت مصرف سهم پردازنده و نرخ پردازش مشاهده‌شده. & نسبت؛ نرخ \\
$\gamma_{si}$ & ضریب خطی $N^{\mathrm{proc}}_{si}/T^{\mathrm{CPU}}_{si}$. & نرخ بر هسته \\
$\mu_{si},U_{si}$ & ظرفیت خطی $y_{si}/u_{si}$ و سقف با حاشیه اطمینان. & درخواست بر ثانیه \\
$\lambda^{\mathrm{obs}}_{si}$ & کل نرخ ورود کانتینر؛ برای پایش و مرجع صف. & درخواست بر ثانیه \\
$R^{\mathrm{plan}}_{si}$ & مجموع تخصیص تازه مبدأهای حل‌شده همین سرویس. & درخواست بر ثانیه \\
$R^{\mathrm{max}}_{esi}$ & ظرفیت عادی مصرف‌نشده برای سطر جاری. & درخواست بر ثانیه \\
$K_{as},\ell_{es}$ & مجموع سقف کانتینرهای محلی و نرخ پذیرفته‌شده مبدأ. & درخواست بر ثانیه \\
$d^{\mathrm{fix}}_{esi},P_s,\delta^Q_{si}$ & هزینه ثابت مسیر، جریمه شکست و مازاد صف. & ثانیه \\
$\kappa_s,r_{esi},w_{esi}$ & ضریب مثبت هزینه، نرخ تخصیص و وزن خروجی. & ضریب؛ نرخ؛ نسبت \\
\end{longtable}
در همه روابط، زمان بر حسب ثانیه و نرخ بر حسب درخواست بر ثانیه است.

\subsubsection{تابع هزینه و مسئله تخصیص}

این بخش فقط حالت $\ell_{es}\leq\sum_iR^{\mathrm{max}}_{esi}$ را حل می‌کند. نرخ نمونه‌های فاقد ظرفیت عادی صفر است و جمع‌های تابع هدف فقط روی $\cI_{es}$ انجام می‌شوند. تابع هزینه مسیر به صورت زیر است:
\begin{equation}
  C^{\mathrm{route}}_{esi}(r)
  =\frac{\kappa_s r}{\mu_{si}-r}
    +d^{\mathrm{fix}}_{esi}r,
  \qquad \kappa_s>0.
  \label{eq:route-cost}
\end{equation}
جمله کسری با نزدیک شدن نرخ تخصیص به ظرفیت کانتینر رشد می‌کند. جمله دوم، هزینه ثابت مسیر را متناسب با بار آن مسیر وارد می‌کند.

با واحدهای جدول و در نظر گرفتن شمار درخواست به عنوان کمیت شمارشی، این تابع شاخص هزینه بی‌واحد است و مستقیما زمان پاسخ یک درخواست نیست. زمان‌های شبکه، خطا و صف پیش از ترکیب، همگی به ثانیه تبدیل می‌شوند.

هدف مدیر خوشه برای زوج $(e,s)$، کمینه کردن مجموع وزن‌دار هزینه‌هاست:
\begin{equation}
  \min_{\mathbf w}\quad
  \mathcal C_{es}(\mathbf w)
  =\sum_{i\in\cI_{es}}
     w_{esi}C^{\mathrm{route}}_{esi}(\ell_{es}w_{esi}),
  \label{eq:weighted-objective}
\end{equation}
با قیود~\eqref{eq:row-weights} و~\eqref{eq:domain}.

برای ساده شدن حل، متغیر نرخ به جای وزن به کار می‌رود. با جایگذاری $w_{esi}=r_{esi}/\ell_{es}$ داریم
\begin{equation}
  \mathcal C_{es}
  =\frac{1}{\ell_{es}}
   \sum_{i\in\cI_{es}}\left[
     \frac{\kappa_s r_{esi}^{2}}
          {\mu_{si}-r_{esi}}
      +d^{\mathrm{fix}}_{esi}r_{esi}^{2}\right].
  \label{eq:objective-substitution}
\end{equation}
ضریب $1/\ell_{es}$ ثابت و مثبت است و حذف آن محل کمینه را تغییر نمی‌دهد. تابع زیر، سهم هر مقصد در هدف تبدیل‌شده را نشان می‌دهد:
\begin{equation}
  \Psi_{esi}(r)
  =\frac{\kappa_s r^2}{\mu_{si}-r}
    +d^{\mathrm{fix}}_{esi}r^2.
  \label{eq:allocation-cost}
\end{equation}
بنابراین، مسئله معادل به صورت زیر نوشته می‌شود:
\begin{equation}
\begin{aligned}
  \min_{\mathbf r}\quad &
       \Phi_{es}(\mathbf r)=\sum_{i\in\cI_{es}}\Psi_{esi}(r_{esi})\\
  \text{با قیود}\quad &
       \sum_{i\in\cI_{es}}r_{esi}=\ell_{es},\\
  &0\leq r_{esi}\leq R^{\mathrm{max}}_{esi}.
\end{aligned}
  \label{eq:allocation-problem}
\end{equation}
توان دوم از ضرب وزن در تابع هزینه و تغییر متغیر حاصل شده است. نرخ ورودی این حل، سهم محلی پذیرفته‌شده است؛ درخواست‌های بیرون از این سهم وارد تابع کسری نمی‌شوند.

\subsubsection{تحدب تابع هدف}

در این بخش، ظرفیت کانتینرها، سقف‌های تخصیص و هزینه‌های ثابت هر حل ثابت‌اند. برای کوتاه شدن روابط، شاخص‌ها حذف می‌شوند و $\mu$، $d$ و $\kappa$ به‌ترتیب ظرفیت کانتینر، هزینه ثابت و ضریب مثبت جمله کسری را نشان می‌دهند.

\Needspace{12\baselineskip}
\begin{proposition}
اگر $\kappa>0$، $\mu>0$ و $d\geq0$ باشند، تابع
$\Psi(r)=\kappa r^2/(\mu-r)+dr^2$
در ناحیه $0\leq r\leq R^{\mathrm{max}}<\mu$ اکیدا محدب است. مشتق دوم این تابع نیز در همین ناحیه اکیدا صعودی است.
\end{proposition}
\begin{proof}
مشتق اول، دوم و سوم تابع عبارت‌اند از
\begin{align}
  \Psi'(r)&=\frac{\kappa r(2\mu-r)}{(\mu-r)^2}+2dr,
    \label{eq:first-derivative}\\
  \Psi''(r)&=\frac{2\kappa\mu^2}{(\mu-r)^3}+2d>0,
    \label{eq:second-derivative}\\
  \Psi'''(r)&=\frac{6\kappa\mu^2}{(\mu-r)^4}>0.
    \label{eq:third-derivative}
\end{align}
در ناحیه مجاز، $\mu-r$ مثبت است. با توجه به $\kappa>0$ و $d\geq0$، رابطه~\eqref{eq:second-derivative} مثبت بودن مشتق دوم و در نتیجه تحدب اکید را نشان می‌دهد. رابطه~\eqref{eq:third-derivative} نیز صعودی بودن مشتق دوم را بیان می‌کند.

مجموع توابع $\Psi_{esi}$ اکیدا محدب است و قیود مسئله~\eqref{eq:allocation-problem} خطی‌اند. ناحیه مجاز، در صورت ناتهی بودن، بسته و کراندار است. بنابراین، حل همین سطر کمینه سراسری یکتا دارد.
\end{proof}

مثبت بودن مشتق دوم، تحدب اکید را ثابت می‌کند؛ مثبت بودن مشتق سوم، صعودی بودن مشتق دوم را نشان می‌دهد. این نتیجه به حل همان سطر با ورودی‌های ثابت مربوط است.

\subsubsection{حل با روش لاگرانژ و نوشتن دستگاه معادلات}
\label{sec:lagrange-system}

برای مبدأ $e$ و میکروسرویس $s$، تعداد نمونه‌های مجاز را $n$ می‌گیریم و در این دستگاه آن‌ها را با شماره‌های $1$ تا $n$ نشان می‌دهیم. تابع هدف وزنی و تابع هدف بر حسب نرخ، از رابطه $\mathcal C_{es}(\mathbf w)=\Phi_{es}(\ell_{es}\mathbf w)/\ell_{es}$ به هم مربوط‌اند.

ابتدا حالتی را می‌نویسیم که جواب در داخل بازه مجاز باشد. قید مجموع وزن‌ها خطی است و مشتق آن نسبت به هر وزن برابر یک است؛ پس برای این حالت می‌توان شرط لاگرانژ را نوشت. مقصدهایی که به سقف می‌رسند، در ادامه جدا بررسی می‌شوند.

\textbf{نوشتن تابع لاگرانژ.}
تابع هدف~\eqref{eq:weighted-objective} را بر حسب وزن‌ها باز می‌کنیم و قید مجموع وزن‌ها را با ضریب $\lambda_{es}\in\mathbb R$ به آن اضافه می‌کنیم:
\begin{equation}
\begin{aligned}
 \mathcal L(\mathbf w,\lambda_{es})={}&
 \sum_{i=1}^{n}\left[
 \frac{\kappa_s\ell_{es}w_{esi}^{2}}
      {\mu_{si}-\ell_{es}w_{esi}}
 +d^{\mathrm{fix}}_{esi}\ell_{es}w_{esi}^{2}\right]\\
 &+\lambda_{es}\left(\sum_{i=1}^{n}w_{esi}-1\right).
\end{aligned}
 \label{eq:lagrangian}
\end{equation}
برای دیدن ارتباط این لاگرانژ با تابع $\Phi_{es}$ در رابطه~\eqref{eq:allocation-problem}، کافی است $\mathbf w=\mathbf r/\ell_{es}$ را جایگذاری کنیم:
\begin{equation}
 \ell_{es}\mathcal L\left(\frac{\mathbf r}{\ell_{es}},\lambda_{es}\right)
 =\Phi_{es}(\mathbf r)
  +\lambda_{es}\left(\sum_{i=1}^{n}r_{esi}-\ell_{es}\right).
 \label{eq:lagrange-scaling}
\end{equation}
سمت راست، تابع هدف بر حسب نرخ و قید مجموع نرخ‌هاست. چون $\ell_{es}$ ثابت و مثبت است، دو نوشتار یک جواب می‌دهند.

\textbf{صفر قرار دادن مشتق‌ها.}
مشتق نسبت به هر وزن و سپس نسبت به $\lambda_{es}$ برابر صفر قرار می‌گیرد:
\begin{equation}
\begin{aligned}
 \frac{\partial\mathcal L}{\partial w_{esi}}
 &={\frac{\kappa_s\ell_{es}w_{esi}
           (2\mu_{si}-\ell_{es}w_{esi})}
          {(\mu_{si}-\ell_{es}w_{esi})^{2}}}
   +2d^{\mathrm{fix}}_{esi}\ell_{es}w_{esi}
   +\lambda_{es}=0,\\
 \frac{\partial\mathcal L}{\partial\lambda_{es}}
 &=\sum_{i=1}^{n}w_{esi}-1=0.
\end{aligned}
 \label{eq:weight-stationarity}
\end{equation}
معادله اول برای تک‌تک مقصدها نوشته می‌شود. این مشتق با جایگذاری $r=\ell_{es}w_{esi}$ در مشتق رابطه~\eqref{eq:first-derivative} به دست می‌آید و سپس $\lambda_{es}$ به آن اضافه می‌شود. معادله دوم نیز مجموع وزن‌ها را یک نگه می‌دارد.

\textbf{تشکیل دستگاه معادلات.}
با جایگذاری $r_{esi}=\ell_{es}w_{esi}$، جمله $\ell_{es}w_{esi}$ در مشتق‌ها به $r_{esi}$ تبدیل می‌شود. قید مجموع نیز به صورت $\sum_i r_{esi}=\ell_{es}\sum_iw_{esi}=\ell_{es}$ درمی‌آید. بنابراین، دستگاه زیر به دست می‌آید:
\begin{equation}
\left\{
\begin{aligned}
 &\frac{\kappa_s r_{es1}(2\mu_{s1}-r_{es1})}
        {(\mu_{s1}-r_{es1})^{2}}
       +2d^{\mathrm{fix}}_{es1}r_{es1}=-\lambda_{es},\\[4pt]
 &\frac{\kappa_s r_{es2}(2\mu_{s2}-r_{es2})}
        {(\mu_{s2}-r_{es2})^{2}}
       +2d^{\mathrm{fix}}_{es2}r_{es2}=-\lambda_{es},\\[2pt]
 &\hspace{4cm}\vdots\\[2pt]
 &\frac{\kappa_s r_{esn}(2\mu_{sn}-r_{esn})}
        {(\mu_{sn}-r_{esn})^{2}}
       +2d^{\mathrm{fix}}_{esn}r_{esn}=-\lambda_{es},\\[4pt]
 &r_{es1}+r_{es2}+\cdots+r_{esn}=\ell_{es}.
\end{aligned}
\right.
 \label{eq:rate-system}
\end{equation}
این دستگاه $n+1$ معادله و $n+1$ مجهول دارد: نرخ‌های $r_{es1}$ تا $r_{esn}$ و ضریب $\lambda_{es}$. ظرفیت‌ها، هزینه‌های مسیر و نرخ تقاضا در این مرحله معلوم‌اند. اگر جواب در داخل بازه مجاز باشد، با توجه به تحدب اکید، نرخ‌های به‌دست‌آمده کمینه یکتای همین مسئله‌اند.

\textbf{بررسی جواب مقید.}
دستگاه~\eqref{eq:rate-system} برای مقصدهایی نوشته شده است که نرخ آن‌ها مثبت و کمتر از سقف باشد. اگر جواب روی مرز قرار بگیرد، برابری مشتق همه مقصدها لازم نیست. برای بیان یکسان این حالت‌ها، مشتق هزینه هر مقصد را با $M_{esi}$ مشخص می‌کنیم:
\begin{equation}
 M_{esi}(r)=
 \frac{\kappa_s r(2\mu_{si}-r)}
      {(\mu_{si}-r)^{2}}
 +2d^{\mathrm{fix}}_{esi}r.
 \label{eq:marginal}
\end{equation}
با تعریف $\nu_{es}=-\lambda_{es}$، جواب معتبر برای مقصدهای دارای سقف مثبت، شرط‌های زیر را رعایت می‌کند:
\begin{equation}
\begin{aligned}
 M_{esi}(r_{esi})&\begin{cases}
 \geq\nu_{es},&r_{esi}=0,\\
 =\nu_{es},&0<r_{esi}<R^{\mathrm{max}}_{esi},\\
 \leq\nu_{es},&r_{esi}=R^{\mathrm{max}}_{esi},
 \end{cases}\\
 \sum_{i\in\cI_{es}}r_{esi}&=\ell_{es},\qquad
 0\leq r_{esi}\leq R^{\mathrm{max}}_{esi}.
\end{aligned}
 \label{eq:optimality-conditions}
\end{equation}
به دلیل تحدب اکید و خطی بودن قیود، هر جواب مجاز با این شرط‌ها، کمینه یکتای همان سطر است.

اگر تقاضا برابر مجموع سقف‌ها باشد، نرخ همه مقصدها از همان سقف‌ها مشخص می‌شود. در این حالت، ممکن است بیش از یک مقدار برای ضریب لاگرانژ معتبر باشد؛ یکتایی نرخ‌ها به معنی یکتایی این ضریب نیست.

روش حل آزاد است؛ جواب باید قیود و شرط‌های~\eqref{eq:optimality-conditions} را با دقت مشخص رعایت کند.

وزن نهایی هر مقصد از نرخ حل‌شده به دست می‌آید:
\begin{equation}
 w_{esi}^{\star}=\frac{r_{esi}^{\star}}{\ell_{es}}.
 \label{eq:final-weights}
\end{equation}
دقت حل و خطای مجاز قید مجموع باید مشخص شوند. پس از گرد کردن وزن‌ها برای اجرا نیز مجموع سطر و سقف مقصدها دوباره بررسی می‌شود؛ تغییر ساده وزن‌ها برای یک شدن مجموع، به‌تنهایی رعایت سقف‌ها را تضمین نمی‌کند.


\subsubsection{اثر پارامترها بر وزن‌ها}
\label{sec:parameter-effects}

اثرهای جدول برای یک سطر با ثابت بودن سایر داده‌ها و سقف‌ها بیان می‌شوند. ظرفیت و سقف هر کانتینر از پایش و تخصیص هسته به دست می‌آیند. با تغییر هر ورودی، وزن‌ها از حل دوباره همان مسئله تعیین می‌شوند.

\Needspace{10\baselineskip}
\begin{longtable}{p{3.2cm}p{10.8cm}}
\caption{اثر پارامترهای کانتینر و مسیر بر تخصیص}\label{tab:local-effects}\\
\toprule پارامتر & اثر تغییر \\
\midrule\endfirsthead
\multicolumn{2}{c}{ادامه جدول \ref{tab:local-effects}}\\
\toprule پارامتر & اثر تغییر \\
\midrule\endhead\bottomrule\endfoot
$\overline d^{\mathrm{net}}_{ei}$ & افزایش پینگ، وزن همان مقصد را کاهش می‌دهد یا ثابت می‌گذارد؛ کاهش پینگ جهت معکوس دارد. \\
$\widehat f_{esi},\delta^Q_{si}$ & با $P_s>0$، افزایش شکست، وزن همان مقصد را کاهش می‌دهد یا ثابت می‌گذارد. افزایش مازاد صف نیز همین اثر را دارد و به مثبت بودن $P_s$ وابسته نیست. \\
$\mu_{si}$ & با افزایش ظرفیت و ثابت بودن هزینه مسیر، فشار هزینه در نرخ ثابت کمتر می‌شود؛ با بازتر شدن سقف، وزن همان مقصد کاهش نمی‌یابد. \\
$R^{\mathrm{max}}_{esi}$ & افزایش سقف، امکان تخصیص بیشتر می‌دهد؛ اگر سقف فعال نباشد، تغییر کوچک آن بی‌اثر است. \\
$g_{si}$ & با ثابت بودن ضریب $\gamma_{si}$، افزایش سهم هسته ظرفیت را خطی بیشتر می‌کند. افزایش سهم تخصیصی باید واقعا در اختیار کانتینر قرار گیرد. \\
$N^{\mathrm{proc}}_{si}$ & با ثابت بودن زمان مصرف و سهم هسته، افزایش شمار پردازش، ظرفیت برآوردشده را بیشتر می‌کند. این اثر، مقایسه جبری ورودی‌هاست و به معنی ثابت ماندن مصرف در هر اجرای واقعی نیست. \\
$T^{\mathrm{CPU}}_{si}$ & با ثابت بودن شمار پردازش و سهم هسته، مصرف بیشتر یعنی هزینه پردازنده بیشتر برای هر درخواست و ظرفیت خطی کمتر. \\
\end{longtable}

\Needspace{10\baselineskip}
\begin{longtable}{p{3.2cm}p{10.8cm}}
\caption{اثر پارامترهای مشترک حل محلی}\label{tab:shared-effects}\\
\toprule پارامتر & نقش \\
\midrule\endfirsthead
\multicolumn{2}{c}{ادامه جدول \ref{tab:shared-effects}}\\
\toprule پارامتر & نقش \\
\midrule\endhead\bottomrule\endfoot
$\ell_{es}$ & با سقف‌های ثابت، افزایش تقاضا نرخ تخصیص مقصدها را کاهش نمی‌دهد؛ وزن نسبی آن‌ها می‌تواند کم یا زیاد شود. \\
$\varepsilon$ & افزایش حاشیه اطمینان، سقف‌ها و ظرفیت قابل پذیرش محلی را کمتر می‌کند. اثر هر وزن به سقف‌های فعال وابسته است. \\
$P_s$ & افزایش آن اهمیت خطا را بیشتر می‌کند؛ جهت تغییر تک‌تک وزن‌ها به نسبت شکست و نرخ همه مقصدها وابسته است. \\
$\kappa_s$ & افزایش این ضریب، اهمیت جمله کسری را نسبت به هزینه ثابت بیشتر می‌کند. جهت تغییر هر وزن عمومی نیست. \\
$\Delta$ & طول پنجره بر تازگی و نوسان اندازه‌گیری اثر دارد. با شمارش و زمان مصرف یک پنجره مشترک، نسبت $N^{\mathrm{proc}}/T^{\mathrm{CPU}}$ به طور مستقیم ضریب جداگانه‌ای از $\Delta$ ندارد. \\
\end{longtable}
نسبت مصرف و نرخ پردازش مستقل از هم نیستند؛ هر دو از همان پنجره ساخته می‌شوند. افزایش مصرف به‌تنهایی نشانه کاهش ظرفیت نیست؛ اگر نرخ پردازش نیز به همان نسبت زیاد شود، برآورد ظرفیت ثابت می‌ماند.

\subsubsection{ثبت تخصیص تازه و پایان محاسبه محلی}

پس از حل سطر $(e,s)$، نرخ‌های تازه در جمع تخصیص همان سرویس ثبت می‌شوند:
\begin{equation}
 R^{\mathrm{plan}}_{si}\gets R^{\mathrm{plan}}_{si}+r_{esi}^{\star}.
 \label{eq:planned-update}
\end{equation}
سطر بعدی از سقف باقی‌مانده همین کانتینر استفاده می‌کند. در پایان حل همه مبدأهای محلی، داریم
\begin{equation}
 R^{\mathrm{plan}}_{si}=\sum_{e\in\mathcal E_a}r_{esi}^{\star}\leq U_{si}.
 \label{eq:aggregate-plan}
\end{equation}
پس جمع نرخ ارسالی همه مبدأها به یک کانتینر از سقف آن بیشتر نمی‌شود. برای سرویس دیگر، جمع تخصیص از صفر آغاز می‌شود و نرخ سرویس قبلی وارد آن نمی‌شود.

ترتیب مبدأها ثابت و ثبت‌شده است. هر سطر با ورودی ثابت، کمینه یکتای خودش را دارد؛ نتیجه کل به ترتیب سطرها وابسته است و کمینه مشترک همه مبدأها ادعا نمی‌شود. جمله کسری، شاخص هزینه تخصیص همین مبدأ است و زمان دقیق صف ناشی از مجموع همه ورودی‌ها نیست.

خروجی این بخش وزن شرطی نمونه‌های محلی است. در ظرفیت کافی، همین وزن‌ها برای ارسال محلی کافی‌اند. برای بخش مازاد تقاضا، سهم خوشه‌های پذیرنده لازم است؛ بخش بعد محاسبه آن را در ریشه توضیح می‌دهد.
